% Part: computability
% Chapter: recursive-functions
% Section: non-pr-functions

\documentclass[../../../include/open-logic-section]{subfiles}

\begin{document}

\olfileid{cmp}{rec}{npr}
\olsection{توابعی که بازگشتی اولیه نیستند}

توابع بازگشتی اولیه همهٔ توابعی را که به‌طور شهودی محاسبه‌پذیرند
دربر نمی‌گیرند. به‌طور شهودی باید روشن باشد که می‌توان فهرستی از همهٔ
توابع یک‌موضعی و بازگشتی اولیه،
$f_0$، $f_1$، $f_2$،~\dots، تهیه کرد، به‌گونه‌ای که بتوانیم مقدار
$f_x$ را برای ورودی $y$ به‌طور مؤثر محاسبه کنیم؛ به بیان دیگر، تابع
$g(x,y)$ که با
\[
g(x,y) = f_x(y)
\]
تعریف می‌شود محاسبه‌پذیر است. اما در این صورت تابع زیر نیز
محاسبه‌پذیر است:
\begin{eqnarray*}
h(x) & = & g(x,x) + 1 \\
& = & f_x(x) +1.
\end{eqnarray*}
به ازای هر تابع بازگشتی اولیهٔ $f_i$، مقدارهای $h$ و~$f_i$ در~$i$
متفاوت‌اند. پس $h$ محاسبه‌پذیر است، اما بازگشتی اولیه نیست؛ دربارهٔ
$g$ نیز می‌توان همین را گفت. این صورت «مؤثر» برهان قطری کانتور است.

می‌توان نمونه‌های صریح‌تری از توابع محاسبه‌پذیری ارائه کرد که بازگشتی
اولیه نیستند. برای نمونه، بگذارید $g^n(x)$ نشان‌دهندهٔ
$g(g(\dots g(x)))$ با مجموعاً $n$ بار کاربرد $g$ باشد، و دنبالهٔ
$g_0,g_1,\dots$ از توابع را چنین تعریف کنید:
\begin{eqnarray*}
g_0(x) & = & x+1 \\
g_{n + 1}(x) & = & g_n^x(x)
\end{eqnarray*}
می‌توانید بررسی کنید که هر تابع $g_n$ بازگشتی اولیه است. هر تابع این
دنباله بسیار سریع‌تر از تابع پیشین رشد می‌کند؛ $g_1(x)$ برابر $2x$
است، $g_2(x)$ برابر $2^x\cdot x$ است و $g_3(x)$ تقریباً مانند برج
نمایی‌ای متشکل از $x$ عددِ~$2$ رشد می‌کند. تابع آکرمان--پِتِر اساساً تابع
$G(x)=g_x(x)$ است و می‌توان نشان داد که سریع‌تر از هر تابع بازگشتی
اولیه رشد می‌کند.

به مسئلهٔ برشمردن توابع بازگشتی اولیه بازگردیم. به یاد آورید که به هر
تابع بازگشتی اولیه نمادگذاری‌ای صوری نسبت داده‌ایم؛ پس کافی است
نمادگذاری‌ها را برشماریم. می‌توانیم به‌صورت بازگشتی به هر نمادگذاری
$F$ عدد طبیعی $\#(F)$ را چنین نسبت دهیم:
\begin{eqnarray*}
\#(0) & = & \langle 0 \rangle \\
\#(S) & = & \langle 1 \rangle \\
\#(\Proj{n}{i}) & = & \langle 2, n, i \rangle \\
\#(\fn{Comp}_{k,l}[H,G_0,\dots,G_{k-1}]) & = & \langle
3,k,l,\#(H),\#(G_0),\dots,\#(G_{k-1}) \rangle \\
\#(\fn{Rec}_l[G,H]) & = & \langle 4, l, \#(G), \#(H) \rangle
\end{eqnarray*}
اینجا از این واقعیت استفاده می‌کنیم که با کدهای بخش پیشین می‌توان هر
دنبالهٔ اعداد را عددی طبیعی در نظر گرفت. نتیجه این است که به هر کد
یک عدد طبیعی نسبت داده می‌شود. البته برخی دنباله‌ها (و بنابراین برخی
عددها) متناظر با هیچ نمادگذاری‌ای نیستند؛ اما می‌توانیم بگذاریم $f_i$
تابع یک‌موضعی و بازگشتی اولیه‌ای باشد که اگر $i$ کد یک نمادگذاری است،
آن نمادگذاری با~$i$ کد شده است؛ و در غیر این صورت تابع ثابت~$0$
باشد. در نتیجه راهی صریح برای برشمردن توابع یک‌موضعی و بازگشتی اولیه
داریم.

(در واقع، برخی توابع مانند تابع ثابت صفر بیش از یک بار در فهرست ظاهر
می‌شوند. این فقط پیامد شیوهٔ کدگذاری ما نیست، بلکه از این واقعیت نیز
ناشی می‌شود که تابع ثابت صفر بیش از یک نمادگذاری دارد. بعداً خواهیم
دید که به‌طور محاسبه‌پذیر نمی‌توان از این تکرارها پرهیز کرد؛ برای
نمونه، هیچ تابع محاسبه‌پذیری نیست که تصمیم بگیرد آیا نمادگذاری
داده‌شده‌ای نمایندهٔ تابع ثابت صفر هست یا نه.)

اکنون می‌توانیم تابع $g(x,y)$ را برابر $f_x(y)$ بگیریم، که در آن
$f_x$ به برشماری‌ای اشاره می‌کند که همین اکنون توصیف کردیم. از کجا
می‌دانیم $g(x,y)$ محاسبه‌پذیر است؟ به‌طور شهودی روشن است: برای
محاسبهٔ $g(x,y)$ نخست $x$ را «بازگشایی» می‌کنیم و می‌بینیم آیا
نمادگذاری تابعی یک‌موضعی است یا نه. اگر چنین باشد، مقدار آن تابع را
برای ورودی~$y$ محاسبه می‌کنیم.

\begin{tagblock}{TMs}
\begin{digress}
شاید از پیش قانع شده باشید که (با قدری کار!) می‌توان برنامه‌ای
(مثلاً در جاوا یا ++C) نوشت که این کار را انجام دهد؛ و اکنون می‌توانیم
به تز چرچ--تورینگ استناد کنیم که می‌گوید هر چیزی که به‌طور شهودی
محاسبه‌پذیر باشد، با ماشین تورینگ قابل محاسبه است.

البته راه مستقیم‌تر برای نشان‌دادن محاسبه‌پذیریِ $g(x,y)$ این است که
ماشین تورینگِ محاسبه‌کنندهٔ آن را به‌صراحت توصیف کنیم. این کار به‌ویژه
از توسل به تز چرچ--تورینگ و شهود پرهیز می‌کند. به‌زودی ابزار کافی
خواهیم ساخت تا با استناد به مدلی از محاسبه که روی ماشین تورینگ
\emph{شبیه‌سازی‌پذیر} است نشان دهیم $g(x,y)$ محاسبه‌پذیر است: یعنی
توابع بازگشتی.
\end{digress}
\end{tagblock}

\end{document}
